%-------------------------%
% Style and macros for the CV. Hand-maintained, /render never rewrites this file.
% Based on Jake Gutierrez's resume template; the custom commands below are Sadra's.
%-------------------------%

\documentclass[letterpaper,11pt]{article}

\usepackage{latexsym}
\usepackage[empty]{fullpage}
\usepackage{titlesec}
\usepackage{marvosym}
\usepackage[usenames,dvipsnames]{xcolor}
\usepackage{verbatim}
\usepackage{enumitem}
\usepackage[hidelinks]{hyperref}
\usepackage{fancyhdr}
\usepackage{lmodern}   % T1 outlines; without it pdflatex renders bitmap fonts
\usepackage[T1]{fontenc}
\usepackage[english]{babel}
\usepackage{tabularx}
\usepackage{xstring}
\usepackage{fontawesome5}
\input{glyphtounicode}

\pagestyle{fancy}
\fancyhf{}
\fancyfoot{}
\renewcommand{\headrulewidth}{0pt}
\renewcommand{\footrulewidth}{0pt}

% Margins
\addtolength{\oddsidemargin}{-0.5in}
\addtolength{\evensidemargin}{-0.5in}
\addtolength{\textwidth}{1in}
\addtolength{\topmargin}{-.5in}
\addtolength{\textheight}{1.0in}

\urlstyle{same}

\raggedbottom
\raggedright
\setlength{\tabcolsep}{0in}

% Section headings
\titleformat{\section}{
  \vspace{-4pt}\scshape\raggedright\large
}{}{0em}{}[\color{black}\titlerule \vspace{-5pt}]

% Machine-readable / ATS-parsable output
\pdfgentounicode=1

%-------------------------%
% Custom commands
\newcommand{\resumeItem}[1]{
  \item\small{
    {#1 \vspace{-2pt}}
  }
}

\newcommand{\resumeSubheading}[4]{
  \vspace{-2pt}\item
    \begin{tabular*}{0.97\textwidth}[t]{l@{\extracolsep{\fill}}r}
      \textbf{#1} & #2 \\
      \textit{\small#3} & \textit{\small #4} \\
    \end{tabular*}\vspace{-7pt}
}

\newcommand{\resumeSubSubheading}[2]{
    \item
    \begin{tabular*}{0.97\textwidth}{l@{\extracolsep{\fill}}r}
      \textit{\small#1} & \textit{\small #2} \\
    \end{tabular*}\vspace{-7pt}
}

\newcommand{\resumeProjectHeading}[2]{
    \item
    \begin{tabular*}{0.97\textwidth}{l@{\extracolsep{\fill}}r}
      \small#1 & #2 \\
    \end{tabular*}\vspace{-7pt}
}

\newcommand{\resumeSubItem}[1]{\resumeItem{#1}\vspace{-4pt}}

\renewcommand\labelitemii{$\vcenter{\hbox{\tiny$\bullet$}}$}

\newcommand{\resumeSubHeadingListStart}{\begin{itemize}[leftmargin=0.01in, label={}]}
\newcommand{\resumeSubHeadingListEnd}{\end{itemize}}
\newcommand{\resumeItemListStart}{\begin{itemize}}
\newcommand{\resumeItemListEnd}{\end{itemize}\vspace{-5pt}}

\definecolor{Black}{RGB}{0, 0, 0}
\newcommand{\seticon}[1]{\textcolor{Black}{\csname #1\endcsname}}

% 1. Define the tag you want to show (the filter)
\def\tagfilter{ALL} % <-- ONLY ENTRIES WITH THIS TAG OR 'ALL' WILL SHOW

% 2. Define the command to display a paper
% Usage: \mypaper{tag1,tag2,...}{Paper Entry Content}
\newcommand{\mytags}[2]{%
    \small
    % Check if the tagfilter is contained within the list of tags (#1)
    % or if 'ALL' is in the list, or if the filter is 'ALL'
    \IfStrEq{\tagfilter}{ALL}{%
        % If tagfilter is ALL, show all papers
        #2\par
    }{%
        % Otherwise, check if the required tag is in the paper's tags
        \IfSubStr{#1}{\tagfilter}{%
            % Tag match found, display the paper content
            #2\par
        }{%
            % No tag match, do nothing (don't display the paper)
        }%
    }%
}

%-------------------------%
% Added for the wiki renderer.

% A wrapping entry: bold title (left) with a date or venue pushed right, then a
% second line for authors or a description. \resumeSubheading's tabular* cannot
% wrap, and publication and project titles routinely overrun the line.
\newcommand{\cventry}[3]{%
  \item\small{%
    \begin{tabularx}{0.97\textwidth}{@{}>{\raggedright\arraybackslash}X@{\hspace{1.2em}}r@{}}
      \textbf{#1} & #2
    \end{tabularx}\\[-3pt]
    #3%
  }\vspace{-3pt}
}

% A dated snapshot of stars/forks/downloads, set under a project's description.
% A PDF has no live badges, so the figures are stamped at render time instead.
\newcommand{\projectstats}[1]{\\{\footnotesize\itshape #1}}

%----------HEADING----------%
\newcommand{\cvheader}{%
  \begin{center}
    \textbf{\Huge \scshape \cvname} \\ \vspace{5pt}
    \seticon{faPhone} \ \small \cvphone \quad
    \href{mailto:\cvemail}{\seticon{faEnvelope} \underline{\cvemail}} \quad
    \href{https://www.linkedin.com/in/\cvlinkedin}{\seticon{faLinkedin} \underline{linkedin.com/\cvlinkedin}} \quad
    \href{https://github.com/\cvgithub}{\seticon{faGithub} \underline{github.com/\cvgithub}}\\ \vspace{2pt}
    \cvlocation
  \end{center}
}
